\section{Changes of height on the toric deformation}\label{sec:local}

Let $Q$ be a lattice polygon with primitive edges in the rank-two
lattice $N$. Translate a vertex to the origin. Write $d_1,\ldots,d_n$
for the oriented edges in cyclic order and $M=\operatorname{Hom}(N,\Z)$.
For Altmann's family $\pi_Q:\mathcal X_Q\to B_Q$ we take the section
$t_1=0$ of simultaneous translation of the edge parameters. Thus
\begin{equation}\label{loc:moments}
 \OO_{B_Q,0}=
 \C\{t_2,\ldots,t_n\}\big/
 \left(\sum_{i=1}^n d_i t_i^k\ :\ k\geq 1\right).
\end{equation}
The first moment identifies the ambient space with $W_Q$, the tangent
space of $B_Q$. For the remaining moments the scheme structure in
\eqref{loc:moments} will matter throughout. The variable $z$ along the
orbits of simultaneous translation is a function on the total space.
We use the convention $\tau_i=z-t_i$, so that $\tau_1=z$.

\subsection{The lift to the total space}
Put $S_\Z=\ker(\Z^n\to N)$, where the $i$th basis vector maps to
$d_i$, and put $\Lambda=\{r\in\Z^n:r|_{S_\Z}=0\}$, using the
coordinate pairing on $\Z^n$. Then $S_\Z^\vee=\Z^n/\Lambda$.
We recall the monomials in Altmann's presentation
\cite[\S\S4.4--5.3]{Alt}. For $c\in M$ choose a vertex of $Q$ at
which $c$ attains its minimum and an edge path from $0$ to that
vertex. If $\lambda_i$ is the signed number of traversals of edge $i$,
the vector with entries $-\lambda_i\langle c,d_i\rangle$ represents
a class $\eta(c)\in S_\Z^\vee$ independent of the choices. A finite
set of monomials $Z_c$ with exponents $(c,\eta(c))$, together with the
$\tau_i$ with exponents $(0,\bar e_i)$, generates the semigroup algebra
used in the construction. Its relations are the equalities of monomials
having the same exponent in $M\oplus S_\Z^\vee$. Base change along the
untranslated moment scheme $B_Q\times\C_z$ gives the total space of
$\pi_Q$ near the central fibre \cite[Theorem~5.1 and \S7]{Alt}.
In particular the presentation uses the subsemigroup generated by these
monomials, even when its semigroup algebra fails to be normal.

Let $T=\operatorname{Hom}(M,\C^*)$ act through the first component of
the exponent. It fixes $z$ and the base. On the open set
$F=\prod_i\tau_i\ne0$ the torus characters can be written
\begin{equation}\label{loc:characters}
             \chi^c=Z_c\prod_i\tau_i^{-\eta_i(c)} .
\end{equation}
For an arbitrary $c$, the notation $Z_c$ denotes the monomial of
exponent $(c,\eta(c))$ in the same algebra
\cite[\S4.6 and \S5.3, Remark~(1)]{Alt}. Here $(\eta_i(c))$ is any
integral representative of $\eta(c)$.
After inverting $F$, all the base monomials are invertible and
the localized semigroup splits off $M$. Indeed the weights $c$ of the
generators generate $M$ as a monoid, since any $c\in M$ occurs in
the dual of the cone over $Q$ at sufficiently large height.
This gives the product with $T$ over the open subset of
$B_Q\times\C_z$ where $F\ne0$.

\begin{proposition}\label{loc:lift}
Let $u$ be coordinates on a smooth germ $D$. Let $H(s,t,u)$ be
holomorphic near the origin, with $H(0,t,u)=0$ and $H_s(0,0,0)\ne0$.
The substitutions
\begin{equation}\label{loc:change}
    t_i'=H(t_i,t,u),\qquad z'=H(z,t,u)
\end{equation}
lift to an automorphism $\Psi_H$ of $\mathcal X_Q\times D$ which fixes
the characters \eqref{loc:characters} on their domain. Its base
automorphism $\varphi_H$ fixes $u$ and preserves the ideal of each
diagonal $t_i=t_j$.
\end{proposition}

\begin{proof}
The moment equations have the following consequence for any convergent
series $f(s)$ whose coefficients depend holomorphically on $t,u,z$:
\begin{equation}\label{loc:series}
              \sum_i r_i f(t_i)=0\qquad(r\in\Lambda).
\end{equation}
Indeed $r$ belongs to the complex row space of the edge matrix, so
every positive power of $s$ gives zero by \eqref{loc:moments}. The
constant term gives zero because $\sum_i d_i=0$. The same calculation
holds in each maximal-ideal quotient, in which it requires only
finitely many powers. Separatedness proves the identity in the analytic
ring. Taking $f(s)=H(s,t,u)^k$ for the two edge rows shows that
\eqref{loc:change} defines a map of $B_Q\times D$.

Set
\[
     U_i=\frac{H(z,t,u)-H(t_i,t,u)}{z-t_i}.
\]
Each divided difference extends holomorphically across $z=t_i$ and
has value $H_s(0,0,0)$ at the origin. Choose a logarithm near this
common value. Apply \eqref{loc:series} to the logarithm of the
divided difference with $s$ in place of $t_i$. Exponentiating gives
\begin{equation}\label{loc:units}
               \prod_i U_i^{r_i}=1 \qquad(r\in\Lambda).
\end{equation}
Thus the assignment $\bar e_i\mapsto U_i$ is a homomorphism from the
whole exponent lattice $S_\Z^\vee$ to the units of the analytic
ring. In particular it respects relations in the saturation of the
integer row span. Multiply each monomial of exponent $(c,\eta)$ by
the value of this homomorphism at $\eta$. On generators this reads
\[
        \tau_i\longmapsto\tau_i U_i,\qquad
        Z_c\longmapsto Z_c\prod_i U_i^{\eta_i(c)} .
\]
Both sides of a binomial acquire the same factor. Since
$\tau_iU_i=z'-t_i'$, these assignments descend to the base change
defining $\mathcal X_Q$ and fix the characters \eqref{loc:characters}.

To invert the construction, lift $H$ to the smooth ambient parameter
space, preserving $H(0,t,u)=0$. There the differential of
$t\mapsto t'$ at the origin is $H_s(0,0,0)$ times the identity in the
$t$ directions. Write $t=b(t',u)$ for its inverse, and let
$J(s',t',u)$ be the inverse in $s$ of $H(s,b(t',u),u)$.
Then $t_i=J(t_i',t',u)$. Formula \eqref{loc:series} for the primed
parameters shows that this inverse descends to $B_Q\times D$.
The divided differences of $J$ pull back to $U_i^{-1}$, so its lift
inverts the map on the total space as well. Finally $t_i'-t_j'$ is
$t_i-t_j$ times a unit, including when $i$ or $j$ equals $1$.
\end{proof}

\subsection{The invariant volume form}
Fix an oriented basis $c_1,c_2$ of $M$. On the open torus family take
the relative form
\begin{equation}\label{loc:form}
       \Omega_Q=dz\wedge d\log\chi^{c_1}\wedge d\log\chi^{c_2}.
\end{equation}
All differentials here and below are relative to $B_Q$ (and to $D$
after taking the product). On families from a fixed Minkowski
decomposition this is Gross's volume form
\cite[Proposition~3.2]{GrossSL}.

\begin{lemma}\label{loc:generator}
The form $\Omega_Q$ generates the relative dualizing line of $\pi_Q$
near the vertex. The automorphism of Proposition~\ref{loc:lift}
satisfies $\Psi_H^*\Omega_Q=H_z(z,t,u)\Omega_Q$, where $H_z$ denotes
the derivative in the height variable.
\end{lemma}

\begin{proof}
Take first a base change to a local Artin $\C$-algebra $A$ with maximal
ideal $\mathfrak m$. The facet opens of the central toric cone cover
the complement of the vertex. For the facet spanned by
$(v_i,1),(v_{i+1},1)$, primitivity of $d_i$ extends this pair to a
lattice basis. The corresponding affine chart is
$\mathbb A^2_{x,y}\times\mathbb G_{m,\zeta}$. In the dual coordinates
the height monomial is $xy$ after absorbing a power of $\zeta$ into
$y$. With a character basis of $T$ adapted to this chart the weights
of $x,y,\zeta$ are $(1,0),(-1,0),(0,1)$.

Lift each facet open by localizing at a semi-invariant lifting its
defining monomial. The resulting affine $A$-algebra is flat, and one
can lift $x,y,\zeta$ in their weight spaces. The lift of $\zeta$ is a
unit. The resulting map from $A[x,y,\zeta^{\pm1}]$ to this algebra
is an isomorphism: on each graded piece for the $\mathfrak m$-adic
filtration it is the identity on the central chart tensored with
$\mathfrak m^k/\mathfrak m^{k+1}$. These algebraic charts apply to
every Artin point of the analytic base, since \eqref{loc:moments} and
the original family are algebraic. In the lifted chart the $T$-invariant
height $z$ is a polynomial $f_i(q)$ in $q=xy$, with
$f_i(q)\equiv q\bmod\mathfrak m$. Its derivative is a unit.
Where $F$ is inverted, $q$ and $x$ are units as well, since $F$
reduces to $q^n$. The characters there are monomials in $x,\zeta$
times invariant units, and the exponents form a unimodular matrix.
The invariant factors have logarithmic differentials proportional to
$dq$, since the invariant ring on this open is $A[q,q^{-1}]$.
Consequently \eqref{loc:form} is
$\pm f_i'(xy)\,dx\wedge dy\wedge d\zeta/\zeta$ on the chart. It
extends as a generator on the punctured fibre over $A$.

Flatness and the Gorenstein property of the central cone give a
relative dualizing line near the vertex, compatible with base change
\cite[Tags 0E4C and 0C08]{Stacks}. On the Artin pullback the local
ring at the vertex has depth three, by flatness and the
Cohen--Macaulay property of the central fibre. The ratio of the form
to a local generator and its reciprocal therefore extend uniquely
across that point.

For descent to the analytic family, choose a local algebraic
generator of the relative dualizing line. In its trivialization
\eqref{loc:form}, which is regular on the algebraic localization at
$F$, has expression $a/F^m$ for some $m$. Multiplication by $F$ is
injective on every Artin pullback: this follows on its associated
graded ring from injectivity of multiplication by $z^n$ on the central
cone. The extension just proved thus gives
$a\in(F^m)+\mathfrak m_{B_Q}^k\OO_{\mathcal X_Q}$ for every $k$.
Krull intersection in the quotient by $(F^m)$ gives $a\in(F^m)$.
The resulting regular section is a generator by its central value.
The formula for $\Psi_H^*\Omega_Q$ follows from
\eqref{loc:form} on inverting $F$. It holds on the family since
injectivity on all the Artin pullbacks and separatedness give
injectivity of multiplication by $F$ in the local ring.
\end{proof}

\subsection{Periods on the link}
Let $L$ be the link with the marking of the central fibre. Over a
local Artin algebra $A$, the deformation is smooth over $A$ in an
annular neighbourhood of $L$. Relative analytic coordinates there
identify the relative de Rham complex locally with the ordinary
holomorphic complex of a polydisc tensored with $A$. The inclusion of
the constant sheaf $A$ is consequently a quasi-isomorphism. Taking
inverse images of sheaves on $L$ gives for any relative three-form
$\omega$ a class
\begin{equation}\label{loc:artin}
                  [\omega]_{L,A}\in H^3(L,\C)\otimes_\C A .
\end{equation}
These identifications commute with maps of Artin algebras. Indeed the
inclusion of constants is natural and a finite good cover computes
constant-sheaf cohomology by a complex of $\C$-vector spaces tensored
with $A$. They also commute with isomorphisms of punctured germs,
using the induced map on central cohomology. One may shrink the annulus
or compute on the punctured germ, since inclusion of the link induces
the same cohomology in either case.

Apply \eqref{loc:artin} to $\Omega_Q$ on the successive maximal-ideal
quotients of $\OO_{B_Q,0}$ to obtain a formal map $\widehat P_1$.
The diagonal cocharacter of the exponent lattice gives $\tau_i$ and
$t_i$ weight one and gives $Z_c$ weight $\sum_i\eta_i(c)$. It fixes
$\chi^c$ and hence gives $\Omega_Q$ weight one. This $\C^*$ acts on
the base, whereas $T$ acts fibrewise. On the punctured central cone
the diagonal cocharacter is homotopic to the identity through torus
actions. Naturality of \eqref{loc:artin} yields
$\widehat P_1(\lambda t)=\lambda\widehat P_1(t)$.
In the completion of a standard graded ring defined by the homogeneous
ideal \eqref{loc:moments}, an element of weight one is its linear
term. Thus $\widehat P_1$ is the restriction of a linear map
\begin{equation}\label{loc:linear}
                  \ell_Q:W_Q\longrightarrow H^3(L,\C).
\end{equation}

Let now $g$ be a holomorphic unit on $\mathcal X_Q\times D$ and let
$h$ be its average over the compact subgroup $(S^1)^2\subset T$.
Take an invariant representative neighbourhood for this average.
Each torus element fixes $\Omega_Q$ and acts trivially on central
link cohomology. Thus $[g\Omega_Q]_{L,A}=[h\Omega_Q]_{L,A}$ for each
Artin point. One can integrate the pullbacks on a compact cycle after
taking smooth de Rham representatives with coefficients in the
finite-dimensional space $A$.

In fact $h$ is holomorphic in $z,t,u$. To see this, lift $g$ to an
ambient convergent series in $Z_c,\tau_i,u$ and take its weight-zero
terms. The exponent monoid of such terms has finitely many generators
by Gordan's lemma. The invariant series converges as a series in
the corresponding monomials: if their maximum total degree is $b$,
a monomial of degree $k$ has at least $k/b$ factors in any
factorization. Choosing one factorization for each term and shrinking
the new polydisc gives absolute convergence by the Cauchy bounds on
the original series. By convexity of $\eta$ \cite[\S4.5(iii)]{Alt},
any product of the $Z_c$ of weight zero has exponent $(0,\bar a)$
for a nonnegative integral vector $a$. It equals $\tau^a$ by the
semigroup relations. Substitution in the finite set of generators
expresses $h$ as a convergent series in $\tau_i,u$, since every
nonconstant generator vanishes at the vertex. Moreover $h$ has the
same nonzero value as $g$ there.

\begin{proof}[Proof of Theorem~\ref{thm:local}]
Choose a convergent representative of $h(z,t,u)$ and apply
Proposition~\ref{loc:lift} to
$H(s,t,u)=\int_0^s h(v,t,u)\,dv$.
The central action of $\Psi_H$ on link cohomology is the identity.
Indeed on the central cone, with $a(z)=H(z,0,0)/z$, its substitutions
are $z\mapsto za(z)$ and $Z_c\mapsto Z_c a(z)^{\sum_i\eta_i(c)}$.
Replacing $a(z)$ by $a(rz)$ for $0\leq r\leq1$ gives an isotopy of
germs to diagonal scaling by the constant $a(0)\ne0$. This scaling
lies in the connected grading torus. By Lemma~\ref{loc:generator} and
the functoriality of \eqref{loc:artin}, the marked class of $g\Omega_Q$
at $(t,u)$ is the class of $\Omega_Q$ at $(t',u)$. On every Artin
point it is therefore $\ell_Q(t')$, where $t'_i=H(t_i,t,u)$.
This expression supplies its convergence. On a disc the usual periods
over transported cycles in a smooth collar are holomorphic functions
whose Taylor coefficients are their relative de Rham classes on the
Artin quotients, so it also gives the usual marked periods there.
Proposition~\ref{prop:variation} identifies \eqref{loc:linear} with
the infinitesimal period isomorphism. The period map consequently
identifies $B_Q$ with its linear image after the base automorphism
$\varphi_H$, uniformly over $D$. The assertion on the marked diagonal
ideals follows from Proposition~\ref{loc:lift}.
\end{proof}
